\chapter{Glossary and Notation}
\label{app:glossary}

\begin{description}[style=standard,itemsep=1pt,parsep=0pt,topsep=3pt]
  \item[AD] Alzheimer's disease.
  \item[ASR] Artifact Subspace Reconstruction, applied by the dataset authors to
  supplied derivatives.
  \item[Balanced accuracy (BA)] Arithmetic mean of per-class recalls; primary
  metric for three-class comparisons.
  \item[Boundary guard] Exclusion of any four-second window intersecting an
  EEGLAB boundary expanded by 0.5 s on each side.
  \item[CAR] Common-average reference across the 19 scalp channels at each time
  sample.
  \item[CN] Cognitively normal control group.
  \item[CP] CANDECOMP/PARAFAC decomposition: sum of tensor rank-one terms.
  \item[FTD] Frontotemporal dementia.
  \item[HOSVD] Higher-order singular value decomposition; used here as a
  training-only Tucker/multilinear projection.
  \item[Inner CV] Cross-validation conducted only inside outer training people
  to choose ranks, features, and classifiers.
  \item[Outer CV] Held-out participant evaluation of the complete selection and
  fitting procedure.
  \item[PSD] Power spectral density.
  \item[Subject pooling] Reduction of multiple windows to one equal-weight
  participant representation.
  \item[Tensor Train (TT)] Chain of low-rank cores whose matrix products recover
  tensor entries.
  \item[Tucker decomposition] Multilinear core tensor multiplied by one factor
  matrix per mode.
  \item[Window] One non-overlapping four-second, 19-channel segment at 250 Hz in
  the processed arrays.
\end{description}

\section{Principal symbols}

\begin{tabular}{ll}
\toprule
\textbf{Symbol} & \textbf{Meaning}\\
\midrule
$\mathcal{X}$ & General multi-way tensor\\
$R$ & CP or terminal latent rank\\
$U,V,G$ & Factor matrices or TT/Tucker cores\\
$A$ & Electrode adjacency matrix\\
$L$ & Electrode graph Laplacian\\
$S$ & Subject delay-covariance moment\\
$b$ & Vectorized channel-lag filter\\
$\omega_i$ & Training subject class weight\\
$M$ & Confusion matrix\\
\bottomrule
\end{tabular}
